\documentclass[10pt,a4paper]{amsart}
\usepackage[margin=19mm]{geometry}
\usepackage{amsmath,amssymb,mathtools}
\usepackage[scr=rsfs,cal=euler]{mathalfa}
\usepackage{xfrac}
\usepackage[unicode,hidelinks,bookmarks=false]{hyperref}
\newcommand{\ie}{i.e.\ }
\newcommand{\cf}{cf.\ }
\newcommand{\resp}{respectively\ }
\newcommand{\Subsection}{\subsection}
\providecommand{\nobreakdash}{\nobreak\hbox{-}}
\setlength{\emergencystretch}{2em}
\multlinegap0pt
\usepackage{amssymb}
\usepackage{mathtools}
\usepackage{algorithm}\usepackage{algpseudocode}
\newtheorem{prop}{Proposition}
\title{Primitive Polynomials of the Form $g(x)+\lambda$ over Finite Fields: Non-Existence Results and Conjectures}
\author{Avnish K. Sharma}
\date{Source: arXiv:2608.07262v2 (CC BY 4.0)}
\begin{document}
\maketitle

\noindent\emph{Source and licence.} Primitive Polynomials of the Form $g(x)+\lambda$ over Finite Fields: Non-Existence Results and Conjectures. Version arXiv:2608.07262v2 (CC BY 4.0), distributed by arXiv under the Creative Commons Attribution 4.0 International licence, http://creativecommons.org/licenses/by/4.0/, as recorded in the arXiv licence metadata for this exact version and shown on the abstract page https://arxiv.org/abs/2608.07262v2. Full licence notice: \texttt{licenses/arxiv-2608.07262-CC-BY-4.0.txt}.\par\smallskip

\noindent\textit{Editorial scope.} Counted unit: the article's prop P4.1 (source0.tex lines 411--414). The article's complete original proof of it is delivered with it (source0.tex lines 415--433, one \texttt{proof} environment of 19 lines). No mathematics is written, repaired or re-derived: the statement and the proof are the article's own lines, in the article's own notation, and no retained line is substituted or re-flowed.  Retained uncounted as the source-local context the proof needs: lines 406--408.  Nothing was omitted from any retained span, so the package carries no omission record.  No retained line contains a citation, so the wrapper carries no bibliography and no bibliography entry.  No retained line contains a figure environment, an image-inclusion command or drawing code, so no image is delivered and the package declares no asset dependency.  Editorial formatting only, in addition: an AMS \texttt{amsart} preamble that restates the article's own declarations and macros used by the retained lines, namely the environments \texttt{prop} on separate counters, together with the standard packages those lines need.  The retained environments print in this document as Proposition 1 in order of appearance, whereas the article prints the same items with the numbering of its own full document.  The complete native source of the article as distributed in the arXiv e-print for this exact version is bundled unchanged as \texttt{sources/207067/source0.tex}.
        \begin{equation}\label{delta}
            \delta:=\sum_{i=0}^{n-1}(-1)^i\lambda^{p^i}\neq 0,
        \end{equation}

        \begin{prop}\label{P4.1}
            Let $p$ be an odd prime. Then, for every $\lambda\in\mathbb{F}_{p^n}$ satisfying
            $\delta\neq 0,$ the polynomial $x^p+x+\lambda$ is irreducible over $\mathbb{F}_{p^n}$ if and only if $n$ is an even positive integer.
        \end{prop}

        \begin{proof}
            
             Suppose $n$ is even. Let $\lambda\in\mathbb{F}_{p^n}$ be such that $\delta\neq 0$. We note that if $\lambda\in\mathbb{F}_p$, then $\lambda^{p^i}=\lambda$ for each $i$, and hence $\delta=\sum_{i=0}^{n-1}(-1)^i\lambda^{p^i}=0$. Therefore, $\lambda\not\in\mathbb{F}_p$. Now, suppose that $\alpha$ is a root of the polynomial $x^p+x+\lambda$. Then
            $$\alpha^p+\alpha+\lambda=0,$$ and hence
            $$\alpha^p=-\alpha-\lambda.$$ 
            If we apply the Frobenius map to $\alpha^p$, then we get 
            $$\alpha^{p^2}=(-\alpha-\lambda)^p=-\alpha^p-\lambda^p=\alpha+\lambda-\lambda^p.$$ 
            Similarly, $$\alpha^{p^3}=(\alpha+\lambda-\lambda^p)^p=\alpha^p+\lambda^p-\lambda^{p^2}=-\alpha-\lambda+\lambda^p-\lambda^{p^2}.$$ 
            Therefore, if $n$ is even then applying the Frobenius map successively to $\alpha^p=-\alpha-\lambda$, we obtain $$\alpha^{p^n}=\alpha+\delta.$$ 
           
            Notice that $\delta^{p^n}=\delta$. Since $\delta\neq 0$, the elements $\alpha$ and $\alpha^{p^n}$ are distinct. Now, we claim that the conjugates of $\alpha$ are given by the following equation. 
            $$\alpha^{{p^n}^j}=\alpha+j\delta \text{  \quad for } j=0,1,2,\ldots,p-1.$$
            Clearly, it is true for $j=0,1$. Now, by induction on $j$, 
            $$\alpha^{{p^n}^j}=(\alpha^{{p^n}^{j-1}})^{p^n}=(\alpha+(j-1)\delta)^{{p^n}}=\alpha^{p^n}+(j-1)\delta^{p^n}=\alpha+\delta+(j-1)\delta=\alpha+j\delta.$$
            Also, notice that if $0\leq r\neq s\leq p-1$, then $\alpha+r\delta\neq \alpha+s\delta$. Hence, all these conjugates are distinct. Since $\alpha$ has $p$ distinct conjugates over $\mathbb{F}_{p^n}$, the polynomial $x^p+x+\lambda$ is irreducible over $\mathbb{F}_{p^n}$.

            Conversely, let the polynomial $x^p+x+\lambda$ is irreducible over $\mathbb{F}_{p^n}$ for some $\lambda$ satisfying Condition \eqref{delta}, and $n$ is odd. Then applying the Frobenius map repeatedly, we obtain $\alpha^{{p^{n}}}=-\alpha-\delta$, and hence $$\alpha^{{p^{2n}}}=(-\alpha-\delta)^{p^n}=-\alpha^{p^n}-\delta^{p^n}=\alpha.$$ 
            This shows that there are at most two distinct conjugates of $\alpha$ over $\mathbb{F}_{p^n}$, and the minimal polynomial of $\alpha$ over $\mathbb{F}_{p^n}$ has degree at most $2$. Thus, the polynomial $x^p+x+\lambda$ is reducible unless $p=2$. Since we have assumed $p$ to be an odd prime, the polynomial $x^p+x+\lambda$ is never irreducible over $\mathbb{F}_{p^n}$.
        \end{proof}

\end{document}
